\chapter{Implementation Source Code}
\label{ch:appendix.b}

If the output of your thesis is computational, you should include properly organized, working source code to ensure that
your reviewer and future readers can reproduce your results. The same practice should, of course, be followed when
authoring any scientific publication with computational aspects.

When submitting your thesis, there are several ways to attach your source code.\footnote{Note that printing it in the
thesis text is \emph{not} one of them.} First, you can package your entire code repository as a ZIP archive and
upload it to the AIS along with your thesis PDF. Alternatively, a more modern approach is to store your code repository
on \emph{\href{https://github.com/}{GitHub}} (or a reliable alternative) and provide a permalink here. This
permalink should ideally point to the specific release or commit version of your code at the time of submission.
The advantage of this method is that you can continue developing your code while still allowing readers to find both
the submitted and the latest versions. Finally, the most robust option is to use an open-science repository,
such as \emph{\href{https://zenodo.org/}{Zenodo}}, where you can upload your code or link your GitHub repository.
Zenodo will archive your code permanently, assign a DOI to it, and generate a Bib\TeX{} reference that you can then
cite here.

Keep in mind that your code repository should remain accessible for the foreseeable future after you submit your thesis.
Because your code is one of your research outputs, it may also be included in your thesis evaluation. Regardless of the
method you choose, always include some documentation with your code, such as a standard \verb|README|
file introducing the repository's structure and providing necessary instructions for running your code. You should then
point your reader to that documentation from this page.
